\chapter{Temporary}

Póki co napiszmy zagadnienia, które były poruszane na wykładzie:
\begin{itemize}
	\item Granica funkcji i ciągłość. Zbieżność po osiach, granice iterowane.
	\item Współrzędne biegunowe. Funkcje rozdzielonych zmiennych.
	\item Klasy funkcji ciągłych.
	\item Parametryzacja krzywej. Poziomica -- idea.
	\item Pochodna kierunkowa (w kierunku wektora) funkcji wielu zmiennych.
	\item Pochodne cząstkowe. Pochodna kierunkowa jako kombinacja liniowa pochodnych cząstkowych. Pochodne złożenia
	\item Gradient i jego własności -- m.in. prostopadłość do poziomicy.
	\item Wektor styczny do krzywej/prosta styczna/płaszczyzna styczna.
	\item Ekstrema lokalne. Punkt przegięcia. Punkt siodłowy. Punkt stacjonarny.
	\item Twierdzenie Hermana Szwarca -- równość pochodnych cząstkowych mieszanych.
	\item Macierz Hesse. Hesjan. Znak hesjanu. Ogólnie: własności i zastosowania.
	\item Schemat wyznaczania ekstremów i punktów siodłowych.
	\item Prosta styczna do poziomicy. Płaszczyzna styczna do poziomicy. Przybliżenia liniowe. Równanie prostej i płaszczyzny stycznej.
	\item Prosta i płaszczyzna styczna do wykresu.
	\item Twierdzenie Lagrange'a (przybliżenie liniowe).
	\item Wyznaczenie płaszczyzny stycznej do wykresy w danym punkcie.
	\item Różniczka zupełna -- metoda przybliżania $f\left( x,y \right) $ przez $df$.
	\item $P\left( x,y \right) dx + Q\left( x,y \right) dy \overset{?}{=} df$ -- forma różniczkowa pierwszego rzędu. Będziemy badać kiedy jest różniczką zupełną.
	\item Punkt krytyczny -- jeśli jest stacjonarny lub któraś pochodna cząstkowa nie istnieje.
	\item Macierz dodatnio określona -- minory etc.
	\item Mnożniki Lagrange'a. 
\end{itemize}
